\problemname{Klippa protein}
Din kompis Björn har nyligen påbörjat ett stort kemiexperiment, med målet att skapa ett hemligt protein. 
Alla protein kan beskrivas som strängar av bokstäver från \texttt{a}-\texttt{z}, exempelvis \texttt{bcaa} och \texttt{hejsan}.
Han behöver skapa ett flertal så-kallade \textit{rena} proteiner. Dessa är protein som bara består av en enda bokstav.
För att göra detta har han mikroskopiska saxar som kallas CRISPR. Med hjälp av dessa kan han ta bort upp till $k$ bokstäver
som ligger bredvid varandra.  

Eftersom han är väldigt lat vill han ha din hjälp att lista ut minsta antalet gånger han behöver klippa ett visst protein för
att det ska bli \textit{rent}.

\section*{Indata}
Den första raden innehåller en sträng som består av bokstäverna \texttt{a}-\texttt{z}, proteinet Björn vill göra rent.
Låt antalet bokstäver i proteinet vara $n$. Då gäller alltid $1 \leq n \leq 32$.

På nästa rad följer ett heltal $k$ ($1 \leq k \leq n$), antalet bokstäver Björn kan ta bort med en klippning. 

\section*{Utdata}
Skriv ut ett heltal: det minsta antalet klippningar Björn behöver göra för att proteinet ska bli rent.


\section*{Poängsättning}
Din lösning kommer att testas på en mängd testfallsgrupper.
För att få poäng för en grupp så måste du klara alla testfall i gruppen.

\noindent
\begin{tabular}{| l | l | p{12cm} |}
  \hline
  \textbf{Grupp} & \textbf{Poäng} & \textbf{Gränser} \\ \hline
  $1$    & $20$       & $k = 1$ \\ \hline
  $2$    & $20$       & Proteinet består endast av bokstäverna \texttt{a} och \texttt{b} \\ \hline
  $3$    & $60$       & Inga ytterligare begränsningar. \\ \hline
\end{tabular}
